\documentclass[10pt,a4paper]{amsart}
\usepackage[margin=19mm]{geometry}
\usepackage{amsmath,amssymb,mathtools}
\usepackage[scr=rsfs,cal=euler]{mathalfa}
\usepackage{xfrac}
\usepackage[unicode,hidelinks,bookmarks=false]{hyperref}
\newcommand{\ie}{i.e.\ }
\newcommand{\cf}{cf.\ }
\newcommand{\resp}{respectively\ }
\newcommand{\Subsection}{\subsection}
\providecommand{\nobreakdash}{\nobreak\hbox{-}}
\setlength{\emergencystretch}{2em}
\multlinegap0pt
\usepackage{tikz}
\usetikzlibrary{cd}
\usepackage{mathtools}
\usepackage{amssymb}
\usepackage{stackengine}
\usepackage{xcolor}
\usepackage{float}
\usepackage{tikz}
\newtheorem{definition}{Definition}
\newtheorem{proposition}{Proposition}
\newcommand{\Post}{\mathrm{Post}}
\newcommand{\Pre}{\mathrm{Pre}}
\newcommand{\Q}{\mathsf{Q}}
\newcommand{\WGPetri}{\mathbf{Petri}^{\mathsf{WG}}}
\newcommand{\profto}{\mathrel{\ooalign{$\longrightarrow$\cr\hfil\!\raisebox{-1.2pt}{\rotatebox{90}{$-$}}\hfil\cr}}}
\title{\bfseries Universal Properties of Petri Net Unfoldings}
\author{Serge Lechenne, Hugo Paquet}
\date{Source: arXiv:2609.17324v1 (CC BY 4.0)}
\begin{document}
\maketitle

\noindent\emph{Source and licence.} \bfseries Universal Properties of Petri Net Unfoldings. Version arXiv:2609.17324v1 (CC BY 4.0), distributed by arXiv under the Creative Commons Attribution 4.0 International licence, http://creativecommons.org/licenses/by/4.0/, as recorded in the arXiv licence metadata for this exact version and shown on the abstract page https://arxiv.org/abs/2609.17324v1. Full licence notice: \texttt{licenses/arxiv-2609.17324-CC-BY-4.0.txt}.\par\smallskip

\noindent\textit{Editorial scope.} Counted unit: the article's proposition prop:Qonmor (source0.tex lines 826--832). The article's complete original proof of it is delivered with it (source0.tex lines 833--866, one \texttt{proof} environment of 34 lines). No mathematics is written, repaired or re-derived: the statement and the proof are the article's own lines, in the article's own notation, and no retained line is substituted or re-flowed.  Retained uncounted as the source-local context the proof needs: lines 655--689.  Nothing was omitted from any retained span, so the package carries no omission record.  No retained line contains a citation, so the wrapper carries no bibliography and no bibliography entry.  The retained lines contain the article's own diagram code, which the wrapper typesets with the standard tikz packages; no image file is delivered and the package declares no asset dependency.  Editorial formatting only, in addition: an AMS \texttt{amsart} preamble that restates the article's own declarations and macros used by the retained lines, namely the environments \texttt{definition}, \texttt{proposition} on separate counters, together with the standard packages those lines need.  The retained environments print in this document as Definition 1, Proposition 1 in order of appearance, whereas the article prints the same items with the numbering of its own full document.  The complete native source of the article as distributed in the arXiv e-print for this exact version is bundled unchanged as \texttt{sources/210587/source0.tex}.
\begin{definition}
  \label{def:morphisms-petri}
  For $P = (S, T, \Pre, \Post, \mu)$ and
  $P' = (S', T', \Pre', \Post', \mu')$, a \emph{map of Petri nets}
  $P \to P'$ consists of a function $\eta : T \rightarrow T'$  and a multirelation $\beta : S \profto S'$ such
  that the diagrams
\[\begin{tikzcd}[row sep=0.2em]
	& S \\
	1 \\
	& {S'}
	\arrow["\beta"{inner sep=.8ex}, "\shortmid"{marking}, from=1-2, to=3-2]
	\arrow["\mu"{inner sep=.8ex}, "\shortmid"{marking}, from=2-1, to=1-2]
	\arrow["{\mu'}"'{inner sep=.8ex}, "\shortmid"{marking}, from=2-1, to=3-2]
\end{tikzcd}
\quad
\begin{tikzcd}
	S & T \\
	{S'} & {T'}
	\arrow["\Pre"{inner sep=.8ex}, "\shortmid"{marking}, from=1-1, to=1-2]
	\arrow["\beta"'{inner sep=.8ex}, "\shortmid"{marking}, from=1-1, to=2-1]
	\arrow["\eta", from=1-2, to=2-2]
	\arrow["{\Pre'}"'{inner sep=.8ex}, "\shortmid"{marking}, from=2-1, to=2-2]
\end{tikzcd}\quad
\begin{tikzcd}
	S & T \\
	{S'} & {T'}
	\arrow["\Post"{inner sep=.8ex}, "\shortmid"{marking}, from=1-1, to=1-2]
	\arrow["\beta"'{inner sep=.8ex}, "\shortmid"{marking}, from=1-1, to=2-1]
	\arrow["\eta", from=1-2, to=2-2]
	\arrow["{\Post'}"'{inner sep=.8ex}, "\shortmid"{marking}, from=2-1, to=2-2]
\end{tikzcd}
\]
commute. (The multisets $\mu$ and $\mu'$ are seen as multirelations from a singleton set, and the function $\eta$ is seen as a multirelation with
$\eta(t, t') = 1$ if $\eta(t) = t'$ and $0$ otherwise.)
\end{definition}

\begin{proposition}
 \label{prop:Qonmor}
  For $P, P' \in \WGPetri$, every rational map $P \xleftarrow{\varphi} R \xrightarrow{\psi} P'$
  induces a morphism of ordinary Petri nets $(\eta, \beta): \Q(P) \to \Q(P')$ where $\eta$ is the
  function $T_R \to T_{P'}$ and $\beta$ is the finitary multirelation
  $\Q(S_P \leftarrow S_R \to S_{P'})$.
\end{proposition}

\begin{proof}
We verify the first condition in  Definition~\ref{def:morphisms-petri} and omit the other two, which use similar arguments. Explicitly, we must show that for every $p' \in S_{P'}$, $\sum_{p \in S_P} \beta(p, p')\mu(p) = \mu'(p')$, where $\mu$ and $\mu'$ are the respective markings of $\Q(P)$ and $\Q(P')$. 

Recall that, by definition of $\Q$, for $p \in S_P$ and $p' \in S_{P'}$, $\beta(p, p')$ is the cardinality of the set $\varphi_S^{-1}\{p\} \cap \psi_S^{-1}\{p'\}$; $\mu(p)$ is the cardinality of the fibre $m_P^{-1}\{ p\}$ for the marking function $m_P : M_P \to S_P$; and similarly $\mu'(p')$ is the cardinality of $m_{P'}^{-1}\{p'\}$. 

By the axioms of rational maps we have the following situation:
\[
\begin{tikzcd}
	S_P & S_R & S_P' \\
	M_P & {M_R} & {M_{P'}} 
	\arrow[from=1-2, to=1-1, "\varphi_S"']
	\arrow[from=1-2, to=1-3, "\psi_S"]
	\arrow[from=2-1, to=1-1, "m_P"]
	\arrow["\lrcorner"{anchor=center, pos=0.125, rotate=180}, draw=none, from=2-2, to=1-1]
	\arrow[from=2-2, to=1-2, "m_R"]
	\arrow[from=2-2, to=2-1]
	\arrow[equals, from=2-2, to=2-3]
	\arrow[from=2-3, to=1-3, "m_{P'}"]
\end{tikzcd}
\]
Thus, by the characterization of pullbacks in $\mathbf{Set}$ we may identify $M_R$ with the coproduct $\sum_{p \in S_P} \varphi_S^{-1}\{p\} \times m_P^{-1}\{p \}$, and under this identification the map $m_R : M_R \to S_R$ is given by $(p, r, j) \mapsto r$. For $p' \in S_{P'}$ we obtain the desired equation by considering the bijection 
\begin{align*}
m_{P'}^{-1}\{p'\} &\cong m_R^{-1} \left( \psi_S^{-1} \{ p' \} \right) \\
&\cong \{ (p, r, j) \mid  p \in S_P, r \in \varphi_S^{-1}\{p \}, j \in m_P^{-1}\{p\} \text{ and } r \in \psi_S^{-1} \{p'\} \} \\ 
&= \{ (p, r, j) \mid  p \in S_P, r \in \varphi_S^{-1}\{p \} \cap \psi_S^{-1}\{p'\}, j \in m_P^{-1}\{p\} \} \\ 
&\cong \coprod_{p \in P} \left(\varphi_S^{-1}\{p \} \cap \psi_S^{-1}\{p'\}\right) \times m_P^{-1}\{p\}
\end{align*}
whose domain has cardinality $\mu(p')$ and whose codomain has cardinality $\sum_{p \in S_P} \beta(p, p') \times \mu(p)$.





\end{proof}

\end{document}
